\pdfoutput=1 % MUST stay within the first 5 lines: tells arXiv to use pdfLaTeX.
%%%%%%%%%%%%%%%%%%%%%%%%%%%%%%%%%%%%%%%%%%%%%%%%%%%%%%%%%%%%%%%%%%%%%%%%%%%%%%%%
%% CPT -- Constraint Projection for Transformers
%% Build:  make        (or: latexmk -pdf main.tex)
%% Submit: make arxiv  -> arxiv-submission.tar.gz, ready to upload
%%
%% EVERY NUMBER IN THIS DOCUMENT COMES FROM AN ARTIFACT IN results/.
%% Each figure and each load-bearing value carries a comment naming its artifact,
%% so any reader can trace it. tests/test_manuscript_consistency.py enforces this.
%%%%%%%%%%%%%%%%%%%%%%%%%%%%%%%%%%%%%%%%%%%%%%%%%%%%%%%%%%%%%%%%%%%%%%%%%%%%%%%%
\documentclass[11pt]{article}

% Options: font=libertinus|times|charter . citestyle=numeric|authoryear
%          paper=letterpaper|a4paper . status={Preprint} . blind . review . nologo
\usepackage[font=libertinus, citestyle=numeric]{kurbanlab}

%%%%%%%%%%%%%%%%%%%%%%%%%%%%%%%%%%%%%%%%%%%%%%%%%%%%%%%%%%%%%%%%%%%%%%%%%%%%%%%%
%% affiliations.tex -- the lab's institution registry.
%%
%% Declare every institution the lab has ever collaborated with, ONCE, here.
%% Then copy this file into every paper. Nothing below appears in a paper unless
%% an \Author actually cites its key, and numbering is assigned in order of first
%% use -- so a solo-HBKU paper shows exactly one affiliation, numbered [1], with
%% no edits to this file.
%%
%% >>> CHECK THESE STRINGS ONCE against how your institutions want to be cited,
%% >>> then never think about them again.
%%%%%%%%%%%%%%%%%%%%%%%%%%%%%%%%%%%%%%%%%%%%%%%%%%%%%%%%%%%%%%%%%%%%%%%%%%%%%%%%

% --- home ---------------------------------------------------------------------


% A second HBKU line, for co-authors who are at HBKU but not in the lab.
\DeclareAffiliation{hbku}{%
  College of Science and Engineering, Hamad Bin Khalifa University, Doha, Qatar}

% --- recurring collaborators --------------------------------------------------
\DeclareAffiliation{tamu}{%
  Department of Computer and Electrical Engineering,
  Texas A\&M University, College Station, TX, USA}

\DeclareAffiliation{iub}{%
  Luddy School of Informatics, Computing, and Engineering,
  Indiana University Bloomington, Bloomington, IN, USA}

\DeclareAffiliation{tamuq}{%
  Department of Electrical and Computer Engineering,
  Texas A\&M University at Qatar, Doha, Qatar}

\DeclareAffiliation{ankara}{%
  Department of Prosthetics and Orthotics,
  Ankara University, Ankara, Turkey}

% --- add new institutions here, never renumber anything -----------------------
% \DeclareAffiliation{qcri}{Qatar Computing Research Institute, Doha, Qatar}
% \DeclareAffiliation{mit} {Department of ..., Massachusetts Institute of Technology, Cambridge, MA, USA}
   % the lab's shared institution registry

%%%%%%%%%%%%%%%%%%%%%%%%%%%%%%%%%%%%%%%%%%%%%%%%%%%%%%%%%%%%%%%%%%%%%%%%%%%%%%%%
%% FRONT MATTER
%%%%%%%%%%%%%%%%%%%%%%%%%%%%%%%%%%%%%%%%%%%%%%%%%%%%%%%%%%%%%%%%%%%%%%%%%%%%%%%%
\title{CPT: Differentiable Constraint Projection for Amortized, Zero-Shot
Retargetable Steering of Autoregressive Generative Models}
\Subtitle{The per-step surrogate targets a concentrated distribution the
sequence-level optimum does not}
\RunningTitle{Differentiable constraint projection for autoregressive models}

%% \Author[<options>]{Name}{affiliation keys, comma separated}
%% Affiliation numbers are assigned in order of first use. Never type a number.
\Author[orcid=0000-0002-1458-302X]{Can Polat}{tamu}
\Author[corresponding=kurbanm@ankara.edu.tr, orcid=0000-0002-7263-0234]{Mustafa Kurban}{tamuq, ankara}
\Author[orcid=0000-0001-9069-770X]{Erchin Serpedin}{tamu}
\Author[corresponding=hkurban@hbku.edu.qa, orcid=0000-0003-3142-2866]{Hasan Kurban}{hbku}

\Keywords{controlled generation; constrained decoding; differentiable optimization
layers; expectation constraints; crystal structure generation}
\CodeURL{https://github.com/KurbanIntelligenceLab/cpt}
\Venue{Preprint}

\begin{document}
\maketitle

\begin{abstract}Generative models are increasingly asked to produce outputs that satisfy a quantitative
requirement rather than merely resemble their training data: a material at a target
density, a molecule within a charge budget, a plan inside a cost. The standard
formalization asks for the distribution closest to the model whose average meets the
target, and that framing has a consequence practitioners rarely confront. An average
says nothing about individual draws, so the cheapest route to it is to spread mass
outward, and the resulting distribution can miss the requirement on most samples anyone
would use. We take the other route, solving a small constrained projection of the next-
token distribution at every decoding step and training through it. The projection is
convex with an analytic derivative, making it a differentiable layer rather than a
decoding heuristic. Two findings follow, on a task small enough to enumerate exactly and
on crystal structure generation from a public materials database. First, the per-step
route is not a worse approximation of the sequence-level solution but an approximation
of a different distribution, and the difference favours it: at equal satisfaction of the
average, the constrained quantity is 6.76 times less variable, and samples violate the
requirement 1.42 times less often than samples from the distribution the field treats as
optimal. Second, because the target enters as a value rather than a training condition,
one trained network serves a family of targets, matching separately trained models at
one seventh of their combined cost. Control of this kind is cheaper than its usual
formalization suggests, and its limits are structural rather than incidental.\end{abstract}

\printkeywords

%%%%%%%%%%%%%%%%%%%%%%%%%%%%%%%%%%%%%%%%%%%%%%%%%%%%%%%%%%%%%%%%%%%%%%%%%%%%%%%%

%%%%%%%%%%%%%%%%%%%%%%%%%%%%%%%%%%%%%%%%%%%%%%%%%%%%%%%%%%%%%%%%%%%%%%%%%%%%%%%%
%% MAIN TEXT. One file per section, in printed order. Figures live in figures/.
%%%%%%%%%%%%%%%%%%%%%%%%%%%%%%%%%%%%%%%%%%%%%%%%%%%%%%%%%%%%%%%%%%%%%%%%%%%%%%%%
\section{Introduction}
\label{sec:intro}

Materials discovery, molecular design and program synthesis increasingly depend on
generative models, and in each the useful output is not merely a plausible one but one that meets a
number: a lattice at a specified density, a composition inside a charge budget, a schedule
within a cost. Learning the distribution of examples does not deliver this, because a model
fit to a corpus reproduces the corpus's spread of the quantity rather than any value a user
picks. Control of a generative model's quantitative behaviour is therefore a prerequisite
for using one as a design tool at all.

The field's standard answer works at the level of whole sequences. Among all distributions
whose expectation of a statistic meets the target, take the one closest in
Kullback--Leibler divergence to the model; the solution is an exponential tilt of the base,
and it is the fixed point that distributional policy gradient and its
descendants pursue. Two properties of that answer are
usually left implicit. It is defined over complete sequences, while an autoregressive model
emits one token at a time, so reaching it requires either sampling from an unnormalized
target or training against an estimator of it.

The second property is the one this paper turns on. An expectation constrains an average,
which says nothing about any individual draw. Constrained optimization has distinguished these
since chance-constrained programming separated expected from per-instance feasibility, and
constrained reinforcement learning inherited the expectation form wholesale. For a generative model the distinction
is not academic, because the samples a user inspects are individual draws rather than
averages, yet the size of the gap between the two notions of feasibility is left
unquantified in this setting.

Working per step is the obvious alternative and has been explored as decoding-time
machinery: grammar and lexical constraints applied as masks, tractable surrogates
consulted during search, and sequential Monte Carlo reweighting partial
sequences. These act on a trained model's output rather than
shaping what the model learns, so a model steered this way was never trained to expect the
steering, and the mismatch between the distribution it learned and the one it is asked to
emit is left for the sampler to absorb.

A separate literature supplies the missing ingredient. Convex programs with implicit
derivatives can sit inside a network as layers, which is what a per-step constraint would
need in order to be trained through rather than bolted on. That
line and the decoding line have not been joined for expectation constraints, so how a
per-step surrogate compares against the sequence-level optimum it approximates has never
been measured against a model actually trained toward that optimum.

Our setting is crystal structure generation, where the target quantity is mass density and
the serialization is text, alongside a synthetic task small enough that every distribution
is computed exactly rather than sampled. Generative models for this domain are established enough that steering one is a fair test
rather than a demonstration; \cref{sec:related} surveys them.

\paragraph{Contributions}

\begin{enumerate}
  \item A differentiable per-step constraint projection with an analytic Jacobian,
        trained through rather than applied at decoding time (\cref{sec:method}).
  \item The concentration result: the surrogate targets a distribution whose
        constrained quantity is concentrated, while the sequence-level optimum
        anti-concentrates it (\cref{sec:concentration}).
  \item Zero-shot retargeting: one checkpoint covers a constraint family with
        statistically bounded equivalence to per-target training
        (\cref{sec:retargeting}).
  \item Four scope boundaries, each measured rather than asserted: a steering
        ceiling, partial and asymmetric two-constraint composition, a serialization
        determinability limit, and inference-time baselines that beat the projection on
        accuracy at a diversity cost (\cref{sec:scope,sec:baselines}).
\end{enumerate}

\begin{kilkey}
The per-step surrogate is not a cheaper approximation of the sequence-level optimum. It
targets a different distribution, one whose constrained quantity is concentrated rather than
spread, and on the tasks measured here that distribution is closer to what a request for a
constraint usually means.
\end{kilkey}

%%%%%%%%%%%%%%%%%%%%%%%%%%%%%%%%%%%%%%%%%%%%%%%%%%%%%%%%%%%%%%%%%%%%%%%%%%%%%%%%

\section{Related work}
\label{sec:related}

\paragraph{Sequence-level distributional control}
Posing controlled generation as a moment constraint and solving for the KL-minimal
distribution yields an exponential-family optimum, and distributional policy gradient
trains a model toward it by importance-weighted sampling against that unnormalized
target~\citep{khalifa2021distributional}. The formulation is attractive because the optimum
is characterized in closed form up to a normalizer and the multiplier has a direct reading
as the price of the constraint. Its practical difficulty is that the target is defined over
complete sequences, so the training signal arrives only once a sequence terminates and the
estimator's variance grows with the tilt strength. Where the resulting distribution is
evaluated, it is evaluated on whether the constrained expectation is met, which is the
quantity the formulation optimizes and not the quantity a user inspecting samples
experiences.

\paragraph{Expectation versus per-instance constraints}
The distinction our concentration result turns on is the one constrained Markov
decision processes make between constraining an expectation and constraining every
trajectory~\citep{altman1999cmdp}, refined by state-wise
formulations~\citep{zhao2023scpo} and by chance-constrained
programming~\citep{charnes1959chance}. The precedent closest in mechanism is the safety
layer, which projects a proposed action onto a feasible set at every
step~\citep{pham2018optlayer}. We adopt that vocabulary: the sequence-level tilt is an
expectation constraint, and what our projection delivers is closer to a per-step
feasibility guarantee whose aggregate happens to meet the expectation.

\paragraph{Per-step constrained decoding}
Lexically constrained decoding enforces predicate-logic constraints during search
without retraining~\citep{lu2021neurologic}, with lookahead heuristics added to reduce
myopia~\citep{lu2022neurologic}. GeLaTo instead distills a tractable generative model to
supply exact conditional guidance~\citep{zhang2023gelato}, and sequential Monte Carlo
methods reweight partial sequences against a target
potential~\citep{zhao2024twisted, loula2025syntactic}. All three act at the per-step locus but treat the
constraint as a filter over candidates or a reweighting of them, so none yields a
distribution that is itself the solution of a constrained problem, and none is
differentiable with respect to the constraint. Ahmed
\etal~\citep{ahmed2023pseudosemantic} are the nearest training-time neighbour, relaxing
a logical constraint into a pseudo-semantic loss over a local approximation, the
same shape of move as reducing a sequence-level expectation to a per-step one, though
approximate in the constraint rather than exact in the support.

\paragraph{Differentiable optimization layers}
Embedding an argmin in a network and backpropagating through its solution is established
for quadratic programs, where differentiating the Karush--Kuhn--Tucker conditions gives the
gradient in closed form~\citep{amos2017optnet}; for general convex programs, where the same
idea is automated by rewriting the problem in a canonical cone
form~\citep{agrawal2019cvxpylayers}; and for hard constraints, where a completion-plus-
correction scheme restores feasibility that a learned map would otherwise
violate~\citep{donti2021dc3}. What these share is a cost profile: a generic solver is
called in the forward pass and its implicit derivative in the backward pass, so the layer
is affordable per example but not per token. None of them has been instantiated as a
divergence projection over a next-token distribution, where the problem is small and
structured enough that a generic solver is the wrong tool.

\paragraph{Crystal structure generation}
Generative models for inorganic crystals include diffusion over invariant
representations~\citep{xie2022cdvae, jiao2023diffcsp}, autoregressive language
models over text serializations~\citep{antunes2024crystallm, gruver2024finetuned}, and
large-scale diffusion trained for property targeting~\citep{zeni2025mattergen}. We use
the same MP-20 subset of the Materials Project~\citep{jain2013materialsproject} and a
text serialization in the spirit of the autoregressive line, because a serialization is
what makes a per-step budget definable at all. That the class is bounded this way is a real
constraint on the method rather than an implementation detail.

%%%%%%%%%%%%%%%%%%%%%%%%%%%%%%%%%%%%%%%%%%%%%%%%%%%%%%%%%%%%%%%%%%%%%%%%%%%%%%%%

\section{Method}
\label{sec:method}

Let $p_\theta(\cdot \mid x_{<t})$ be the model's next-token distribution over a
vocabulary $\mathcal{V}$, let $m_t \in \{0,1\}^{|\mathcal{V}|}$ be a mask marking tokens
the grammar permits after the prefix $x_{<t}$, and let $\phi \in \R^{|\mathcal{V}|}$
assign each token an increment of the constrained quantity. Given a target total $c$
for a sequence of length $L$, define the running budget
\begin{equation}
  b_t = \frac{c - \sum_{s<t} \phi_{x_s}}{L - t + 1},
  \label{eq:budget}
\end{equation}
the per-step increment that would land the total on target if every remaining step
matched it. The layer solves
\begin{equation}
  q_t = \argmin_{q \in \Delta(\mathcal{V})}
        \KL\!\left(q \;\|\; p_\theta(\cdot \mid x_{<t})\right)
  \quad \text{subject to} \quad
  \inner{q}{\phi} = b_t, \quad \operatorname{supp}(q) \subseteq \operatorname{supp}(m_t).
  \label{eq:projection}
\end{equation}
This is a KL projection onto the intersection of an affine slice with a face of the
simplex. Its solution is an exponential tilt of the masked model distribution,
$q_t \propto m_t \odot p_\theta \exp(\lambda \phi)$, with the multiplier $\lambda$ fixed
by the constraint; a one-dimensional monotone root find gives it to machine precision.
\Cref{si:solver} verifies the solve and its gradient oracle. Two properties matter
downstream. The support restriction is satisfied
\emph{exactly}, because masked tokens receive zero mass by construction. The
expectation constraint is satisfied exactly \emph{when it is attainable}: if $b_t$ lies
outside the range of $\phi$ over the allowed tokens, no distribution on that support
meets it, and the layer falls back to the unconstrained masked distribution and records
the fallback.

\paragraph{Differentiating the solve}
The Jacobian of $q_t$ with respect to the incoming logits follows from the implicit
function theorem applied to the scalar stationarity condition, and reduces to a rank-one
correction of a diagonal, so the backward pass costs a few vector operations rather
than a linear solve. Because the correction involves only inner products with $\phi$, the
Hessian-vector products needed for the multiplier's sensitivity can be batched; we
measure a speedup from single precision of
%% results/A2_training_summary.json (hessian_speedup_vs_batch)
$1.33\times$ to $2.64\times$ over the batch sizes tested (\cref{si:overhead}).

\paragraph{What the mask does and does not do}
On the crystal task the mask is a character-level grammar over a text serialization. It
removes the syntactic failure mode completely, since every sampled string parses, but it is
not a validity oracle; \cref{si:parse} separates the two effects.

%%%%%%%%%%%%%%%%%%%%%%%%%%%%%%%%%%%%%%%%%%%%%%%%%%%%%%%%%%%%%%%%%%%%%%%%%%%%%%%%

\section{Experiments and results}
\label{sec:experiments}

\subsection{Shared setup}
\label{sec:setup}

Two settings carry every result below, and both are described here once rather than
repeated per experiment. The \emph{synthetic} setting is a length-8 sequence over a 4-symbol vocabulary, so all
%% results/WS2_synthetic_ground_truth.json (task.n_sequences, task.regime_S_feasible_count)
$65{,}536$ sequences enumerate and every distribution is computed rather than sampled: no
divergence reported on this task is an estimate. A per-symbol multiplicity budget defines
the support constraint, satisfied by $36{,}120$ sequences, and each symbol carries a
charge; the expectation constraint asks for a mean total charge of zero. The \emph{crystal}
setting is the MP-20 subset of the Materials Project, serialized as text under the grammar
of \cref{sec:method}, where the constrained quantity is mass density in
\unit{\gram\per\centi\metre\cubed}.

Four training arms recur throughout, differing only in what they train through rather than
in architecture or mask. The \emph{mask-only} arm trains with the grammar mask and no
projection. The \emph{projection} arm trains through \cref{eq:projection} and is read
either without the layer, giving its backbone, or with it, giving the deployed system. The
\emph{sequence-level comparator} trains by maximum likelihood on exact samples from the
sequence-level tilt, which is distributional policy gradient's fixed point with unit
importance weights; that rule is chosen deliberately, because it gives the comparator the
best case of the method it stands for, so no gradient-estimator artefact can be blamed for
the outcome. Baselines needing no training are named where they are used.

Metrics follow the constraint. On the synthetic task we report the mean of the constrained
quantity, its variance, and the per-sequence violation rate, meaning the probability mass
on sequences whose own total misses the target. On the crystal task we report absolute
density error, structural validity under an independent parser, and unique-formula fraction
as a diversity measure. Every figure reported below is read from a stored artifact at build
time rather than transcribed.

\subsection{The surrogate targets a concentrated distribution}
\label{sec:concentration}

\paragraph{The comparator reaches its target}
%% results/WS7_a6_tilt_comparator.json (summary.kl_from_tilt, summary.mean_charge)
The comparator lands $0.0167 \pm 0.0048$ nats from the exact tilt with a mean charge of
$-0.062 \pm 0.122$ against a target of zero and no mass outside the mask. It is not a
failed run, which is what makes the comparison in \cref{tab:concentration} informative
rather than a win over something broken.

\begin{table}[t]
  \centering
  %% Every value: results/A2_training_summary.json (arms.*, [mean, sd] PAIRS -- index [0])
  %% and results/WS7_a6_tilt_comparator.json (summary.*, reference.tilt,
  %% reference.base_S_conditioned). Exact enumeration, no sampling.
  \caption{\textbf{Both arms satisfy the expectation constraint; only one concentrates
    it.} The sequence-level optimum's charge variance exceeds the unconstrained data's,
    while the projection's is far below it. Per-sequence violation is the mass on
    sequences whose own total charge is nonzero. Mean $\pm$ s.d.\ over seeds where an
    arm is trained; the tilt and the data distribution are exact.}
  \label{tab:concentration}
  \small
  \begin{tabular*}{\textwidth}{@{\extracolsep{\fill}}l S[table-format=+1.4] S[table-format=1.4] S[table-format=2.3]@{}}
    \toprule
    \kilth{Arm} & {\kilth{Mean charge}} & {\kilth{Violation}} & {\kilth{Variance}} \\
    \midrule
    Data, mask-conditioned (exact)   & +2.5164 & 0.9109 &  8.564 \\
    Mask only                        & +2.4694 & 0.9083 &  8.762 \\
    Sequence-level comparator        & -0.0619 & 0.8751 & 10.362 \\
    Tilt, sequence-level optimum     & +0.0000 & 0.8755 & 10.452 \\
    \midrule
    Projection (deployed)            & \best{+0.0170} & \best{0.6155} & \best{1.547} \\
    \bottomrule
  \end{tabular*}
\end{table}

\begin{figure}[t]
  \centering
  %% figures/fig1_charge_distributions.pdf <- results/WS7_charge_marginals.json
  %% Exact enumerated distributions marginalized onto integer total charge.
  \kilgraphics{figures/fig1_charge_distributions.pdf}
  \caption{\textbf{The optimum spreads the constrained quantity; the projection
    concentrates it.} Exact charge distributions over all sequences. The
    sequence-level optimum and the model trained toward it place mass at both extremes
    to buy a mean of zero, ending wider than the mask-only arm they started from. The
    projection instead spikes at the target. Legend entries give each arm's charge
    variance.}
  \label{fig:charge}
\end{figure}

\paragraph{Result}
Both the projection and the comparator satisfy the expectation constraint, and the comparator's seed spread on mean charge
is wider than the gap between them, so we claim no ordering there. The difference that
lies outside seed noise is in the last two columns. The
optimum's charge variance exceeds the mask-conditioned data's, \emph{raising} the variance of
the quantity being constrained by $22.1\%$, while the projection lowers it by $81.9\%$. Against the optimum rather than against the unconstrained arm, the
projection's variance is lower by a factor of $6.76$, a gap exceeding three of the comparator's seed
standard deviations. Individual samples follow: the optimum's own mass sits on
constraint-violating sequences $1.42$ times more often than the projection's does.

The mechanism is visible in \cref{fig:charge} and follows from what the tilt is. It is
the KL-closest distribution with the right mean, and reaching a mean of zero from a base
centred at $+2.469$ is cheapest by moving mass symmetrically outward. Nothing in the
objective rewards concentration. Bounding an expected increment at every step, by
contrast, bounds the total in expectation and narrows it as a side effect, because each
step's feasible set is centred on what the remaining budget can still absorb.
%% results/WS1_induced_vs_tilt.json (comparison.excess_kl_of_induced_over_tilt_nats)
The surrogate's $1.025$ nats of excess divergence from the optimum is therefore a trade
rather than a loss, and \cref{si:frontier} reports how that overhead varies with
constraint strength.

\subsection{One checkpoint retargets across the constraint family}
\label{sec:retargeting}

Because the constraint enters \cref{eq:budget} as a target value rather than as a
training condition, a trained checkpoint can be deployed at values it never trained for.
We test this on the crystal task, where the constrained quantity is mass density and the
%% results/WS4_zero_shot_retargeting.json (matrix: 3 checkpoints x 10 targets)
%% results/WS4_retargeting_seed_replication.json (summary.*)
family is ten density targets spanning the data distribution. Three checkpoints, each
trained at one target, are deployed across all ten. Energies quoted here come from a
surrogate (\cref{si:energy}), not density functional theory.

Structural validity is exact at every one of the thirty cells, including targets no
checkpoint trained for and targets beyond the reachable range. The mask is a property of the
grammar, not of the training target, so retargeting cannot break it.

On accuracy, one checkpoint deployed across the sub-ceiling family is level with the
frontier of
%% preregistration/WS4_equivalence_margin.md fixes the margin BEFORE the test;
%% results/WS4_equivalence.json reads it and fails if the file disagrees.
$7$ separately trained models. A difference test is the wrong instrument here. A large
$p$-value is absence of evidence, not evidence of equivalence, and at $n = 7$ paired
targets the smallest attainable two-sided Wilcoxon $p$ is $2/2^7 = 0.0156$ under any
data, which \cref{si:floor} verifies by exhaustive enumeration. We therefore pre-registered an equivalence margin, defined as the retargeted
arm's own mean within-target seed standard deviation, on the reasoning that a difference
below the method's own reproducibility is not actionable and that a margin tied to the
arm's own noise cannot be widened to pass without changing the arm. Two one-sided tests
against that margin of $0.189$ give $p = 0.0023$, and the pre-registered secondary
condition ($p < 0.05$ at a $0.10$ margin) also holds (\cref{si:equivalence}).

The honest form of the claim is parity, not superiority: on a single seed the retargeted
checkpoint appeared \emph{better} than the frontier, but replication across seeds found
it better at only $3$ of $7$ targets. What survives is that one training run is
statistically indistinguishable from seven at one seventh of the training compute, which
is a claim about cost structure rather than about accuracy. \Cref{si:tab:perseed} gives the
per-seed values behind both readings.

\subsection{Where the method stops}
\label{sec:scope}

\paragraph{Steering saturates at a backbone ceiling}
Requested and achieved density track closely at low targets and then flatten
%% figures/fig2_steering_frontier.pdf <- results/FIG2_steering_frontier.json,
%% aggregated seed-wise from results/WS3_steering_frontier_rows.json
(\cref{fig:frontier}). The plateau sits at $10.405$~\unit{\gram\per\centi\metre\cubed}:
requests beyond it are not met, and no amount of additional tilt reaches them, because
the projection can only redistribute mass the backbone already assigns. The constraint
is a reweighting of the model's own proposals, not a source of new modes.

\begin{figure}[t]
  \centering
  \begin{subfigure}[t]{0.48\linewidth}
    %% figures/fig2_steering_frontier.pdf <- results/FIG2_steering_frontier.json
    \kilgraphics{figures/fig2_steering_frontier.pdf}
    \caption{Requested against achieved density, with the ceiling marked.}
    \label{fig:frontier}
  \end{subfigure}\hfill
  \begin{subfigure}[t]{0.48\linewidth}
    %% figures/fig4_infeasibility_vs_target.pdf <- results/WS3_steering_frontier_rows.json
    %% (infeas_init / infeas) and scripts/crystal/analyze_sweep.py's repair factor
    \kilgraphics{figures/fig4_infeasibility_vs_target.pdf}
    \caption{Infeasible steps before and after training, with the repair factor.}
    \label{fig:infeas}
  \end{subfigure}
  \caption{\textbf{The ceiling and its mechanism.} Steering follows the request until the
    backbone runs out of mass to redistribute; the repair the layer performs weakens as the
    request hardens, and where it approaches unity the ceiling binds.}
\end{figure}

\Cref{fig:infeas} gives the mechanism. Before training, a large fraction of decoding
steps present a budget outside the attainable range, and the layer falls back. Training
repairs most of that at easy targets and progressively less at hard ones; where the
repair factor approaches unity, the ceiling binds.

\paragraph{Two constraints compose partially, and asymmetrically}
The second additive constraint is valence electron density, chosen because it is
computable for every element in the split and only partly correlated with mass density
%% results/CP7_second_channel.json (channel_statistics.correlation = 0.725)
($r = 0.725$). The joint solve converges but steers weakly. Across
%% results/WS5_two_constraint_grid.json (summary.*, fixed_valence_columns)
$7$ trained cells, the fraction of the requested valence shift that survives to the
output never exceeds $11.3\%$, and the second constraint's effect is dominated by
budget exhaustion: at most positions the mass budget the first constraint has already
committed leaves nothing for the second to allocate.

The asymmetry is sharper than the aggregate. Asked to \emph{lower} valence density the
constraint moves the right way in $2$ of $2$ cells; asked to \emph{raise} it, in $2$ of
$5$. A running budget is a remaining quantity: spending less of it is always available by
choosing lighter or fewer atoms, while spending more than the prefix has left is not
available once the budget is exhausted. \Cref{si:grid} gives every cell against
its own control. Charge neutrality, the constraint one would most want here, is not
expressible in this class at all: it quantifies existentially over oxidation states,
and collapsing to one charge per element destroys the predicate rather than approximating
it (\cref{si:expressible}).

\paragraph{The constraint class depends on the serialization}
A per-step budget requires that the prefix determine it. Mass density is a ratio, so a
%% results/CP9_serialization_expressibility.json (rows, verdict; 393 structures)
target density fixes a mass budget only once the cell volume is known. Over $393$
structures and $1{,}217$ element positions, a serialization that closes the cell
parameters first makes the budget determinable at $100\%$ of positions, while one that
emits composition first makes it determinable at $0\%$. The constraint class is
therefore a property of the serialization, not of the layer, and a target quantity that
no prefix determines cannot be steered by this mechanism at all.

\subsection{Baselines}
\label{sec:baselines}

\paragraph{Conditional training is more accurate, where it has data}
Training a single model conditioned on the target value is the strongest same-cost
baseline. It is more accurate than the projection wherever the training split has
%% results/WS3_conditional_baseline.json (summary.mae_below_mean, novel_formula_fraction)
support: mean absolute error $0.085$ against $0.196$ below the data mean. It pays for that
in coverage, emitting $5.9\%$ novel formulas against $94.1\%$, and its accuracy degrades
where the conditioning slice thins, precisely where a practitioner asking for an unusual
target is operating.

\paragraph{Best-of-$n$ beats the projection on accuracy, and pays in diversity}
Given enough inference compute, drawing $n$ samples from the unsteered model and keeping
%% results/WS6_best_of_n.json (rows for error, diversity block for unique fraction)
the closest one needs no training and wins on accuracy: it crosses the projection's error
at $n = 16$ (\cref{fig:pareto}). The cost is concentration of the accepted set. Unique
formula fraction falls from $0.824$ at $n=1$ to $0.228$ at $n=64$, while the
projection's $0.782$ does not depend on how hard the target is to hit. Sequential Monte
Carlo, the principled member of the same family, reaches the projection at
%% results/WS6_smc_steering.json (rows; sigma_sweep for the untuned caveat)
$6$ of $10$ targets within $K=64$ particles but is worse than best-of-$n$ at equal
compute by a mean factor of $4.2$; its potential width is untuned, so that figure is a
lower bound on what SMC could achieve rather than a refutation of it.

\begin{figure}[t]
  \centering
  %% figures/fig3_accuracy_diversity_pareto.pdf <- results/FIG3_accuracy_diversity_pareto.json,
  %% built from results/WS6_best_of_n.json and results/WS3_conditional_baseline.json
  \kilgraphics[0.62\linewidth]{figures/fig3_accuracy_diversity_pareto.pdf}
  \caption{\textbf{Accuracy bought by selection is diversity spent.} Mean best-error
    against unique-formula fraction over the sub-ceiling targets. Best-of-$n$ walks down
    and to the left as $n$ grows, trading output variety for accuracy; the projection sits
    off the curve at one forward pass per output.}
  \label{fig:pareto}
\end{figure}

The honest reading is that the projection's advantage is not accuracy per target. It is
that the cost is paid once at training and never again at inference, while the
inference-time family pays per output forever and concentrates its outputs to do so.

%%%%%%%%%%%%%%%%%%%%%%%%%%%%%%%%%%%%%%%%%%%%%%%%%%%%%%%%%%%%%%%%%%%%%%%%%%%%%%%%

\section{Discussion}
\label{sec:discussion}

The result we did not expect is that the sequence-level optimum anti-concentrates the
quantity it constrains. Stated after the fact it is almost obvious, since an expectation
constraint says nothing about individual draws and the cheapest way to move a mean is to move
the tails. It nonetheless inverts the usual reading of a per-step surrogate as a compromise. The distinction between constraining an expectation and constraining every
trajectory is old; what we add is a measurement of how far apart those
two objects are for a concrete generative model, and the observation that the per-step route
lands closer to the second.

The boundaries are as informative as the results. The steering ceiling says the
projection reweights rather than creates. The composition asymmetry says a running budget
is a remaining quantity, so constraints that want more of it fail differently from
constraints that want less. The determinability result says the expressible constraint
class is chosen when the serialization is chosen, which is a design decision usually made
for other reasons. And best-of-$n$'s accuracy win says the contribution is cost structure,
not accuracy: a claim weaker than the one we set out to make and better supported than it
would have been. \Cref{si:defects} records three errors caught during
this work and the mechanism that caught each, since each would otherwise have produced a
plausible-looking wrong number.

%%%%%%%%%%%%%%%%%%%%%%%%%%%%%%%%%%%%%%%%%%%%%%%%%%%%%%%%%%%%%%%%%%%%%%%%%%%%%%%%

\section{Conclusion}
\label{sec:conclusion}

Solving a small differentiable constrained projection at every decoding step, and training
through it, produces a model whose samples individually respect a quantitative target far
more often than samples from the distribution the field treats as optimal for that target.
That inversion is the paper's substantive finding, and it has a mechanism rather than a
coincidence behind it: an expectation constraint is indifferent to individual draws, so the
cheapest route to a mean is through the tails, while a per-step feasible set is centred on
what the remaining budget can absorb. Because the target enters as a value rather than a
training condition, the same trained network serves an entire family of targets, which
relocates the cost of control from inference to a single training run.

\paragraph{Limitations}
The projection reweights the backbone's own proposals, so it cannot reach a target the
backbone assigns no mass to, and the achievable range saturates well before the requested
one does. Two correlated constraints sharing a single token budget compose only partially,
and asymmetrically: the direction that spends less of the budget is reliably available
while the direction that spends more is not. The expressible class is fixed by the
serialization, since a target no prefix determines cannot be steered at all, and charge
neutrality falls outside that class rather than merely being unimplemented. On accuracy per
target the method is not the strongest option available, because sampling repeatedly from
an unsteered model and selecting the closest output wins given enough inference compute.
The synthetic results are exact but small, and the crystal energetics rest on a surrogate
rather than density functional theory.

\paragraph{Future work}
Three threads follow directly. The two-constraint grid is one row short of the range needed
to separate a smooth dependence on target compatibility from a jump, and completing it is a
matter of compute rather than method. The particle-filter comparison used an untuned
potential width, so the figure reported here bounds that family from below rather than
characterising it, and a tuned version would sharpen the comparison the paper draws. The
serialization result is currently a statement about two orderings measured on one quantity;
turning it into a predictive account of which constraints a given serialization admits needs
several serializations budgeted together, which is the form in which it would become a
design rule rather than an observation.

%%%%%%%%%%%%%%%%%%%%%%%%%%%%%%%%%%%%%%%%%%%%%%%%%%%%%%%%%%%%%%%%%%%%%%%%%%%%%%%%


%% BACK MATTER
%%%%%%%%%%%%%%%%%%%%%%%%%%%%%%%%%%%%%%%%%%%%%%%%%%%%%%%%%%%%%%%%%%%%%%%%%%%%%%%%
\begin{kilrepro}
Every number here is read from a JSON artifact in \texttt{results/} at build time, never
transcribed, and each figure and load-bearing value carries a comment naming its artifact;
the test suite recomputes them and fails if the prose disagrees.
\end{kilrepro}

\begin{contributions}
C.P.\ designed the layer, ran the experiments, and drafted the paper. H.K.\ supervised the
project and revised the manuscript.
\end{contributions}

\begin{acknowledgments}
We thank the members of the Kurban Intelligence Lab for discussions.
\end{acknowledgments}

\begin{availability}
Code and artifacts:
\ifbool{KIL@blind}
  {\url{https://anonymous.4open.science/r/XXXXXX}}
  {\url{https://github.com/KurbanIntelligenceLab/cpt}}.
The MP-20 subset used here is public.
\end{availability}

\begin{conflicts}
The authors declare no competing interests.
\end{conflicts}

%%%%%%%%%%%%%%%%%%%%%%%%%%%%%%%%%%%%%%%%%%%%%%%%%%%%%%%%%%%%%%%%%%%%%%%%%%%%%%%%
\bibliography{references}

%%%%%%%%%%%%%%%%%%%%%%%%%%%%%%%%%%%%%%%%%%%%%%%%%%%%%%%%%%%%%%%%%%%%%%%%%%%%%%%%

%%%%%%%%%%%%%%%%%%%%%%%%%%%%%%%%%%%%%%%%%%%%%%%%%%%%%%%%%%%%%%%%%%%%%%%%%%%%%%%%
%% SUPPLEMENTARY INFORMATION. Single file; \appendix switches the numbering to S.
%%%%%%%%%%%%%%%%%%%%%%%%%%%%%%%%%%%%%%%%%%%%%%%%%%%%%%%%%%%%%%%%%%%%%%%%%%%%%%%%
\appendix
\renewcommand{\thesection}{S\arabic{section}}
\renewcommand{\thefigure}{S\arabic{figure}}
\renewcommand{\thetable}{S\arabic{table}}
\setcounter{section}{0}\setcounter{figure}{0}\setcounter{table}{0}

\section{Solver verification and the gradient oracle}
\label{si:solver}

The projection defined in the method is solved by a one-dimensional monotone root
find on the multiplier. Two things are verified rather than assumed. First, the solve
%% results/CP2_hessian_speedup_sweep.json, results/CP6_multiconstraint.json
converges to the constraint to machine precision wherever the target is attainable, and
reports non-convergence rather than returning a near-miss where it is not; the
non-convergence flag is what the fallback path keys on. Second, the analytic Jacobian
agrees with finite differences and with autograd through an unrolled solver.

For the multi-constraint form, convergence alone is insufficient: a converged solve on a
near-singular system is not trustworthy. We therefore report the conditioning of the
%% results/CP6_multiconstraint.json (conditioning report, vocabulary sweep)
constraint matrix alongside every joint solve, and sweep vocabulary size rather than
testing a single configuration, because a single grid at one vocabulary would have read as
``constraints always compose''.

\section{Overhead and timing}
\label{si:overhead}

%% results/A2_training_summary.json (wall_clock_ratio_A3_over_A2, hessian_speedup_vs_batch)
Training through the projection costs $6.0\times$ the wall clock of training with the
mask alone, on the synthetic task at matched steps. The backward pass is a rank-one
correction rather than a linear solve, and batching the Hessian-vector products in
single precision gives a speedup ranging from $1.33\times$ at the smallest batch to
$2.64\times$ at the largest. At inference the layer adds one root find per step.

\section{Grammar mask and parse validity}
\label{si:parse}

%% results/WS2_parse_validity.json (arms.masked, mask_freedom, claim_scope)
The character-level grammar makes every sampled string parse: the syntactic parse rate
is $1.00$ under an independent parser. It is not a validity oracle: only $47\%$ of
those strings build into a structure, the remainder failing on semantic grounds such as
overlapping sites or unphysical stoichiometries, which parse without complaint. The mask
leaves real freedom: a mean of $4.14$ allowed characters per step, though $70.6\%$ of
steps are fully forced by the grammar.

\section{How the overhead varies with constraint strength}
\label{si:frontier}

%% results/WS1_steering_frontier.json (summary.*)
The surrogate's excess divergence over the sequence-level optimum is close to fixed in
absolute terms while the optimum's own cost varies over a $129$-fold range as the
constraint pulls harder. The relative waste therefore falls from $97.8\%$ at the weakest
steering to $32.4\%$ at the strongest: the surrogate is least efficient exactly where the
constraint barely matters, and most efficient where it binds. \Cref{si:solver} covers the
solve itself and \cref{si:overhead} its wall-clock cost.

\section{Energy screen calibration}
\label{si:energy}

%% results/WS3_energy_calibration.json (surrogate, held_out_agreement)
Formation energies are estimated with CHGNet~\citep{deng2023chgnet} at
version v0.3.0~\citep{chgnet_v030} rather than density functional theory. On $1{,}456$ held-out structures the surrogate agrees with its reference at
$0.0210$~\unit{\electronvolt\per atom} mean absolute error and Pearson $r = 0.9996$.
This is a surrogate screen and is labelled as such wherever it is used; no claim in the
main text depends on a DFT-quality energy.

\section{Statistical floor of the difference test}
\label{si:floor}

%% results/WS6_best_of_n.json and the enumeration in scripts/crystal/equivalence_test.py
The claim that a paired Wilcoxon test at $n = 7$ cannot return $p < 0.0156$ is verified
by exhaustive enumeration over all $2^{7}$ sign patterns rather than by appeal to a
table. This is why the equivalence procedure, rather than the
difference test, carries the retargeting claim.

\section{The full equivalence appendix}
\label{si:equivalence}

%% preregistration/WS4_equivalence_margin.md (the binding margin and decision rule)
%% results/WS4_equivalence.json (tost_p, bootstrap interval, margin_matches_stated_rule)
The margin is defined before the test, in a pre-registration file that the analysis
script reads; the script fails if the file is absent or if its value disagrees with the
stated rule. The rule is that the margin equals the retargeted arm's mean within-target
seed standard deviation. Two one-sided tests give $p = 0.0023$ at that margin of
$0.189$, and the secondary pre-registered condition, still $p < 0.05$ at a $0.10$
margin, also holds. A bootstrap interval is reported alongside the parametric one and
the wider of the two is what the main text quotes.

%% results/WS4_retargeting_seed_replication.json (summary.*)
The difference test is reported for completeness and interpreted as uninformative: the
paired Wilcoxon over $7$ targets returns $p = 1.0$, but the smallest attainable
two-sided $p$ at that sample size is $2/2^{7} = 0.0156$ under any data whatsoever, so a
large value carries almost no information.

%% results/WS4_retargeting_seed_replication.json (per_seed, summary)
The table below gives the three replication seeds behind both readings, and it is
included because the aggregate hides the reason to distrust one of them. Each row is a
separate training run of the projection arm at the same target, read first without the layer
to give its backbone and then with the layer across the whole sub-ceiling family. The
backbone column is what makes seed~0 unusable as evidence: at
$22.769$ it sits far above the other two, which land at
$8.381$ and $8.692$, so its lower
sub-ceiling error of $0.4406$ reflects a backbone that happened
to propose high-density structures rather than a layer that steered better. Averaging the
three gives $0.6497$ against the frontier's
$0.6620$, which is the comparison the main text reports, and
the spread of $0.1506$ across seeds is what the pre-registered
margin is derived from. All three seeds exceed the published ceiling of
$10.405$, so the excursion is a matter of degree rather than of kind.

\begin{table}[t]
  \centering
  %% results/WS4_retargeting_seed_replication.json and results/WS4_zero_shot_retargeting.json
  \caption{\textbf{Per-seed retargeting summary.} Densities in
    \unit{\gram\per\centi\metre\cubed}. Seed 0's backbone is an excursion, which is why
    its lower error is not the reading to take; the published ceiling is $10.405$.
    Structural validity is $1.0$ in every cell of every seed.}
  \label{si:tab:perseed}
  \small
  \begin{tabular*}{\textwidth}{@{\extracolsep{\fill}}@{}l S[table-format=2.3] S[table-format=2.3] S[table-format=1.4]@{}}
    \toprule
    \kilth{Seed} & {\kilth{Backbone mean}} & {\kilth{Max achieved}} &
    {\kilth{Sub-ceiling MAE}} \\
    \midrule
    0 & 22.769 & 15.495 & 0.4406 \\
    1 &  8.381 & 10.439 & 0.7889 \\
    2 &  8.692 & 10.596 & 0.7197 \\
    \midrule
    Mean of seeds        & {--} & {--} & 0.6497 \\
    Frontier, 7 runs     & {--} & {--} & 0.6620 \\
    \bottomrule
  \end{tabular*}
\end{table}

\section{Two-constraint grid, per cell}
\label{si:grid}

%% results/WS5_two_constraint_grid.json (summary.fixed_valence_columns)
\Cref{si:tab:grid} gives every cell of the two-constraint grid at a shared valence
target, each against its own single-constraint control. The control is necessary: the
achieved valence density rises with the density target even when no valence constraint
is applied, so an absolute level would show apparent steering where there is none. Only
the deviation from control is the second constraint's effect.

\begin{table}[t]
  \centering
  %% results/WS5_two_constraint_grid.json (summary.fixed_valence_columns[*].rows)
  \caption{\textbf{Direction of the second constraint's effect.} Cells at a shared
    valence target across density rows. ``Asked'' is the direction the target sits
    relative to that cell's control. Lowering the quantity succeeds in both cells that
    request it; raising it succeeds in two of five.}
  \label{si:tab:grid}
  \small
  \begin{tabular*}{\textwidth}{@{\extracolsep{\fill}}@{}S[table-format=1.3] S[table-format=2.1] S[table-format=1.4]
                    S[table-format=1.4] S[table-format=+1.4] l c@{}}
    \toprule
    {\kilth{Valence}} & {\kilth{Density}} & {\kilth{Achieved}} & {\kilth{Control}} &
    {\kilth{Deviation}} & \kilth{Asked} & \kilth{Correct} \\
    \midrule
    0.282 &  3.0 & 0.2447 & 0.2424 & +0.0023 & up   & yes \\
    0.282 &  5.0 & 0.3062 & 0.3141 & -0.0079 & down & yes \\
    0.282 &  8.0 & 0.3251 & 0.3956 & -0.0705 & down & yes \\
    0.484 &  3.0 & 0.2603 & 0.2424 & +0.0179 & up   & yes \\
    0.484 &  8.0 & 0.3482 & 0.3956 & -0.0474 & up   & no  \\
    0.686 &  3.0 & 0.2353 & 0.2424 & -0.0071 & up   & no  \\
    0.686 &  8.0 & 0.3707 & 0.3956 & -0.0248 & up   & no  \\
    \bottomrule
  \end{tabular*}
\end{table}

\section{The boundary of the expressible class}
\label{si:expressible}

%% results/CP7_second_channel.json (charge_is_not_linear, channel_statistics)
%% results/CP9_serialization_expressibility.json (rows, verdict)
Charge neutrality cannot be written as a per-step linear expectation over this
vocabulary. The screen that defines it quantifies existentially over oxidation-state
assignments; most elements in the split are multivalent; and collapsing to a single
charge per element destroys the predicate rather than approximating it. This is a
property of the constraint, not a gap in the implementation, and it bounds the class the
layer addresses. Valence electron density was adopted in its place because it is
computable for every element present and only partly correlated with mass density
($r = 0.725$), so a joint solve is not one constraint written twice.

A second boundary is the serialization. Over $393$ structures and $1{,}217$ element
positions, a cell-first serialization makes the density budget determinable at $100\%$
of positions and a composition-first serialization at $0\%$.

\section{Defects found, and how}
\label{si:defects}

Three classes of error were caught during this work by mechanisms worth reporting, since
each would have produced a plausible-looking wrong number.

A budget generalization initially read cell-parameter token payloads as physical values
when they are bin indices, which inflated volumes while still producing budgets that
looked reasonable; it was caught by an equivalence check against the original
single-constraint path rather than by inspection. A charge-variance comparison was made
against the unconstrained arm when the claim required the sequence-level optimum:
a true statement answering a different question, caught in review, and the corrected
figure is stronger than the original. A crossing count in an earlier draft was a literal
in the document-generating script rather than a value computed from the data, and
survived a verification pass that checked every number in the same table; it is now
recomputed and asserted by a test.

The general lesson we would pass on is that verifying numbers is not the same as
verifying prose. A test that recomputes a table's values will not catch the sentence
describing them.


\end{document}
